%%%%%%%%%%%%%%%%%%%%%%%%%%%%%%%%%%%%%%%%%%%%%%%%%%%%%%%%%%%%%%%%%%%%%%%%%%%%%%%
%%  TKS_Psychology_Weaponized_Therapy.tex
%%  Comprehensive TKS Formalization for Psychological Profile Analysis
%%  and Intervention Design Systems
%%
%%  Part of TKS v7.4 - December 2025
%%  ADDITIVE to v7.4 canon
%%%%%%%%%%%%%%%%%%%%%%%%%%%%%%%%%%%%%%%%%%%%%%%%%%%%%%%%%%%%%%%%%%%%%%%%%%%%%%%

\documentclass[11pt,twoside]{book}

%% ============================================================================
%% PACKAGES
%% ============================================================================
\usepackage[utf8]{inputenc}
\usepackage[T1]{fontenc}
\usepackage[margin=1in,inner=1.2in,outer=1in]{geometry}
\usepackage{amsmath,amssymb,amsthm}
\usepackage{mathtools}
\usepackage{stmaryrd}
\usepackage{enumitem}
\usepackage{booktabs}
\usepackage{longtable}
\usepackage{ltablex}
\usepackage{tabularx}
\usepackage{array}
\usepackage{tcolorbox}
\tcbuselibrary{breakable,skins}
\usepackage{tikz-cd}
\usepackage{xcolor}
\usepackage{hyperref}
\usepackage{ragged2e}

%% ============================================================================
%% FORMATTING SETTINGS
%% ============================================================================
\setlength{\parindent}{0pt}
\setlength{\parskip}{0.5em}
\renewcommand{\arraystretch}{1.2}

%% ============================================================================
%% CUSTOM COMMANDS (from TKS canonical)
%% ============================================================================
\newcommand{\noe}[1]{\nu_{#1}}
\newcommand{\Ideas}{\mathcal{I}}
\newcommand{\Worlds}{\mathcal{W}}
\newcommand{\Elements}{\mathcal{E}}
\newcommand{\Noetics}{\mathcal{N}}
\newcommand{\Foundations}{\mathcal{F}}
\newcommand{\Acquisitions}{\mathcal{A}}
\newcommand{\Frac}{\Phi}
\newcommand{\TransFrac}[1]{\Phi^{#1}}
\newcommand{\MultiFrac}[2]{\Phi^{#1}_{#2}}
\newcommand{\RPM}{\mathsf{RPM}}
\newcommand{\ACBE}{\mathsf{ACBE}}
\newcommand{\HausdorffDim}{\dim_H}
\newcommand{\FractalSpectrum}{\sigma_\Phi}
\newcommand{\fix}{\mathrm{fix}}
\newcommand{\sem}[1]{\llbracket #1 \rrbracket}
\newcommand{\Tadd}{\oplus_T}
\newcommand{\Tsub}{\ominus_T}
\newcommand{\Tmul}{\otimes_T}
\newcommand{\Tdiv}{\oslash_T}
\newcommand{\Wadd}{\oplus_W}
\newcommand{\Wsub}{\ominus_W}
\newcommand{\Failed}{\mathtt{Failed}}

%% New commands for this document
\newcommand{\Lacunarity}{\Lambda}
\newcommand{\PsychAttractor}{\mathcal{A}_\psi}
\newcommand{\ManipFrac}{\Phi_{\mathrm{int}}}
\newcommand{\ProfileFrac}{\Phi_{\mathrm{prof}}}
\newcommand{\VulnSet}{\mathcal{V}}
\newcommand{\StealthFactor}{\sigma_{\mathrm{stealth}}}
\newcommand{\TextCorpus}{\mathcal{T}}
\newcommand{\NoetMap}{\mathcal{M}_\nu}
\newcommand{\IFS}{\mathcal{I}_S}

%% Theorems
\newtheorem{definition}{Definition}[chapter]
\newtheorem{theorem}{Theorem}[chapter]
\newtheorem{proposition}{Proposition}[chapter]
\newtheorem{lemma}{Lemma}[chapter]
\newtheorem{corollary}{Corollary}[chapter]
\newtheorem{example}{Example}[chapter]
\newtheorem{remark}{Remark}[chapter]

%% ============================================================================
%% DOCUMENT
%% ============================================================================
\begin{document}

\title{\textbf{Psychology: Weaponized Therapy} \\[0.5em]
\Large TKS Mathematical Framework for Psychological Profile Analysis \\
and Targeted Intervention Design}
\author{TKS v7.4 Formalization Project}
\date{December 2025}
\maketitle

\tableofcontents

%% ============================================================================
%% PART I: THEORETICAL FOUNDATIONS
%% ============================================================================
\part{Theoretical Foundations}

\chapter{Mathematical Framework for Psychological Profiling}
\label{ch:psychological-profiling}

\begin{tcolorbox}[colback=blue!5!white,colframe=blue!75!black,title=\textbf{Purpose and Scope},breakable]
This document establishes the mathematical foundations for psychological profile extraction,
attractor computation, vulnerability detection, and intervention design using the TKS canonical
framework. The system provides rigorous formalization of:
\begin{enumerate}
    \item \texttt{extract\_fractal\_from\_text()}: Text-to-noetic-fractal mapping
    \item \texttt{psych\_attractor = attractor(profile)}: IFS fixed-point computation
    \item \texttt{vulnerabilities = find\_high\_lacunarity\_regions()}: Gap detection
    \item \texttt{manipulation\_fractal = design\_fractal()}: Counter-fractal engineering
\end{enumerate}
The mathematical treatment is objective and descriptive, documenting the system's operational
principles without normative evaluation.
\end{tcolorbox}

%% ============================================================================
%% SECTION: FRACTAL EXTRACTION FROM TEXT
%% ============================================================================
\section{Profile Extraction: \texttt{extract\_fractal\_from\_text()}}

\begin{definition}[Text Corpus]
\label{def:text-corpus}
A \emph{text corpus} $\TextCorpus = \{t_1, t_2, \ldots, t_n\}$ is a collection of textual
utterances produced by subject $S$, including written communications, speech transcripts,
social media posts, and behavioral descriptions.
\end{definition}

\begin{definition}[Noetic Mapping Function]
\label{def:noetic-mapping}
The \emph{noetic mapping function} $\NoetMap : \TextCorpus \to \Frac$ is defined by:
\begin{enumerate}[label=(\roman*)]
    \item \textbf{Lexical Analysis}: Parse $\TextCorpus$ into semantic units $\{u_1, \ldots, u_m\}$
    \item \textbf{Noetic Classification}: Map each unit to dominant noetic:
    \[
        \kappa : u_i \mapsto k_i \in \{0, 1, \ldots, 9\}
    \]
    \item \textbf{Sequence Construction}: Form the fractal:
    \[
        \ProfileFrac = \langle \kappa(u_1) : \kappa(u_2) : \cdots : \kappa(u_m) \rangle
    \]
\end{enumerate}
\end{definition}

\begin{definition}[Noetic Classification Rules]
\label{def:noetic-class-rules}
The classification function $\kappa$ maps semantic content to noetics as follows:

\begin{center}
\footnotesize
\renewcommand{\arraystretch}{1.2}
\begin{tabularx}{\textwidth}{c|l|X}
\toprule
\textbf{Noetic} & \textbf{Name} & \textbf{Textual Indicators} \\
\midrule
$\noe{0}$ & Idea & Pure statements, definitions, abstractions \\
$\noe{1}$ & Mind & Awareness, attention, ``I notice'', ``I realize'' \\
$\noe{2}$ & Positive & Attraction, desire, ``I want'', ``I like'', approach \\
$\noe{3}$ & Negative & Rejection, avoidance, ``I hate'', ``I fear'', withdrawal \\
$\noe{4}$ & Vibration & Energy, intensity, excitement, emphasis, activation \\
$\noe{5}$ & Female & Reception, absorption, ``I absorb'', passivity, integration \\
$\noe{6}$ & Male & Projection, action, ``I do'', ``I create'', expression \\
$\noe{7}$ & Rhythm & Patterns, habits, cycles, ``always'', ``every time'' \\
$\noe{8}$ & Above & Causation, origin, ``because'', ``the reason is'' \\
$\noe{9}$ & Below & Effects, results, ``therefore'', ``it leads to'' \\
\bottomrule
\end{tabularx}
\end{center}
\end{definition}

\begin{theorem}[Profile Extraction Canonical Form]
\label{thm:profile-extraction}
For any text corpus $\TextCorpus$ with sufficient data ($|\TextCorpus| > \theta_{\min}$),
the extracted profile fractal converges to a characteristic signature:
\begin{equation}
\boxed{
    \texttt{extract\_fractal\_from\_text}(\TextCorpus) =
    \lim_{n \to \infty} \left\langle \kappa(u_1) : \cdots : \kappa(u_n) \right\rangle
    = \ProfileFrac_S
}
\end{equation}
where $\ProfileFrac_S$ is the \emph{psychological fingerprint} of subject $S$.
\end{theorem}

\begin{proof}
By the ergodic property of noetic sequences (TKS v7.3, Theorem 6.41), the empirical
distribution of noetics in $\TextCorpus$ converges to the subject's characteristic
distribution $\pi_S : \{0, \ldots, 9\} \to [0,1]$. The fractal $\ProfileFrac_S$ encodes
this distribution in its sequence statistics.
\end{proof}

%% ============================================================================
%% PLAIN ENGLISH: EXTRACT FRACTAL
%% ============================================================================
\begin{tcolorbox}[colback=green!5!white,colframe=green!65!black,title=\textbf{Plain English: The ``Mind-Reading Machine''},breakable]

\textbf{The Simple Version:}
\texttt{extract\_fractal\_from\_text()} is like a ``psychological X-ray'' that reads someone's
writing and speech to decode their mental operating system. Just as an X-ray reveals bone
structure beneath skin, this function reveals psychological structure beneath words.

\textbf{The Fingerprint Scanner Metaphor:}
Every person has a unique fingerprint---a pattern of ridges and whorls. Similarly, every
person has a unique \textit{psychological fingerprint}---a pattern of how they think, feel,
and process reality. This function scans their ``mental fingerprint'' from their words.

\begin{itemize}
    \item \textbf{The scanner} = The \texttt{extract\_fractal\_from\_text()} algorithm
    \item \textbf{The finger pressed on glass} = The text corpus (their words, posts, emails)
    \item \textbf{The fingerprint image} = The extracted fractal $\ProfileFrac_S$
    \item \textbf{Ridge patterns} = Noetic sequences (how their mind moves between states)
    \item \textbf{Unique whorls} = Their characteristic noetic distribution
\end{itemize}

\textbf{How It Works (Step by Step):}
\begin{enumerate}
    \item \textbf{Collect samples}: Gather their writing, speech, social media posts
    \item \textbf{Parse into units}: Break text into meaningful chunks
    \item \textbf{Classify each chunk}: ``Is this attraction? Rejection? Causation?''
    \item \textbf{Build the sequence}: Chain the noetics in order
    \item \textbf{Extract the pattern}: Find the characteristic fingerprint
\end{enumerate}

\textbf{What It Reveals:}
\begin{itemize}
    \item How often they approach vs.\ avoid (ratio of $\noe{2}$ to $\noe{3}$)
    \item Whether they think in causes or effects ($\noe{8}$ vs.\ $\noe{9}$)
    \item Their energy patterns (presence of $\noe{4}$)
    \item Whether they act or receive ($\noe{6}$ vs.\ $\noe{5}$)
    \item Their habitual thought cycles ($\noe{7}$ patterns)
\end{itemize}

\end{tcolorbox}

%% ============================================================================
%% SECTION: PSYCHOLOGICAL ATTRACTOR COMPUTATION
%% ============================================================================
\section{Attractor Computation: \texttt{psych\_attractor = attractor(profile)}}

\begin{definition}[Induced Iterated Function System]
\label{def:induced-ifs}
Given a profile fractal $\ProfileFrac_S = \langle k_1 : k_2 : \cdots : k_n \rangle$,
the \emph{induced IFS} is:
\[
    \IFS(\ProfileFrac_S) = \{(\noe{k_i}, r_{k_i}, p_{k_i}) : i = 1, \ldots, n\}
\]
where:
\begin{itemize}
    \item $\noe{k_i} : \Ideas \to \Ideas$ is the noetic contraction mapping
    \item $r_{k_i} \in (0, 1)$ is the contraction ratio
    \item $p_{k_i} \in (0, 1)$ is the probability weight (from frequency in $\ProfileFrac_S$)
\end{itemize}
\end{definition}

\begin{theorem}[Psychological Attractor Existence]
\label{thm:psych-attractor-existence}
For any profile $\ProfileFrac_S$ with non-trivial noetic content, there exists a unique
psychological attractor $\PsychAttractor_S \subset \Ideas$ satisfying:

\begin{enumerate}[label=(\roman*)]
    \item \textbf{Invariance}: $\bigcup_{k \in S} \noe{k}(\PsychAttractor_S) = \PsychAttractor_S$
    \item \textbf{Attraction}: $\lim_{n \to \infty} d_H(\ProfileFrac^n_S(X), \PsychAttractor_S) = 0$
          for any compact $X \subset \Ideas$
    \item \textbf{Minimality}: No proper closed subset satisfies (i) and (ii)
\end{enumerate}
\end{theorem}

\begin{proof}
By the Hutchinson-Barnsley theorem (TKS v7.3, Theorem 5.12), any IFS with contractive
maps on a complete metric space has a unique attractor. The noetic operators satisfy
the contraction condition with ratios $r_k < 1$ under the noetic metric $d_\nu$.
\end{proof}

\begin{definition}[Attractor Computation Algorithm]
\label{def:attractor-algorithm}
The psychological attractor is computed via Kleene iteration:
\begin{equation}
\boxed{
    \texttt{attractor}(\ProfileFrac_S) = \fix\left(\lambda X. \bigcup_{k \in S} \noe{k}(X)\right)
    = \bigsqcup_{n=0}^{\infty} F^n(\bot)
}
\end{equation}
where $F(X) = \bigcup_{k \in S} \noe{k}(X)$ and $\bot$ is the bottom element of the domain.
\end{definition}

\begin{proposition}[Attractor Properties]
\label{prop:attractor-properties}
The psychological attractor $\PsychAttractor_S$ encodes:
\begin{enumerate}[label=(\roman*)]
    \item \textbf{Hausdorff Dimension}: $\HausdorffDim(\PsychAttractor_S) = \frac{\log|S|}{\log(1/\bar{r})}$
          where $\bar{r}$ is the geometric mean contraction ratio
    \item \textbf{Support of Stable States}: The set of psychological configurations the
          subject gravitates toward
    \item \textbf{Self-Similarity}: The attractor is self-similar at all scales of analysis
\end{enumerate}
\end{proposition}

%% ============================================================================
%% PLAIN ENGLISH: ATTRACTOR COMPUTATION
%% ============================================================================
\begin{tcolorbox}[colback=green!5!white,colframe=green!65!black,title=\textbf{Plain English: The ``Psychological Gravity Map''},breakable]

\textbf{The Simple Version:}
\texttt{attractor(profile)} computes where someone's mind ``naturally falls''---their
psychological ``lowest energy state.'' Just as a marble in a bowl always rolls to the
bottom, a mind always gravitates toward its attractor.

\textbf{The Gravity Well Metaphor:}
Imagine space-time curved by a massive object---things naturally fall toward the center
of the gravity well. The psychological attractor is where all mental states ``fall toward.''
It's the person's psychological center of gravity.

\begin{itemize}
    \item \textbf{The gravity well} = The attractor $\PsychAttractor_S$
    \item \textbf{Objects falling in} = Mental states converging over time
    \item \textbf{The shape of the well} = The attractor's structure
    \item \textbf{Multiple wells} = Multiple attractors (complex personalities)
    \item \textbf{The depth} = How strongly they're pulled to that state
\end{itemize}

\textbf{What It Computes:}
\begin{enumerate}
    \item Their ``default'' mental/emotional configuration
    \item What states they always return to after perturbation
    \item The stable patterns in their psychology
    \item How complex their psychological structure is (dimension)
    \item The ``shape'' of their mental landscape
\end{enumerate}

\textbf{Types of Attractors Found:}
\begin{itemize}
    \item \textbf{Fixed Point}: Single stable state (always returns to same mood)
    \item \textbf{Limit Cycle}: Oscillates between states (mood swings)
    \item \textbf{Strange Attractor}: Complex, chaotic but bounded (unpredictable but within limits)
\end{itemize}

\end{tcolorbox}

%% ============================================================================
%% SECTION: VULNERABILITY DETECTION
%% ============================================================================
\section{Vulnerability Detection: \texttt{find\_high\_lacunarity\_regions()}}

\begin{definition}[TKS Lacunarity (Canonical)]
\label{def:tks-lacunarity-vuln}
The \emph{lacunarity} of a psychological fractal $\Frac$ at scale $\varepsilon$ is:
\[
    \Lacunarity(\Frac, \varepsilon) := \frac{\mathrm{Var}[M_\varepsilon]}{\mathbb{E}[M_\varepsilon]^2} + 1
\]
where $M_\varepsilon$ measures the occupancy of $\varepsilon$-boxes covering the attractor.
\end{definition}

\begin{definition}[Vulnerability Region]
\label{def:vulnerability-region}
A \emph{vulnerability region} $\VulnSet \subset \PsychAttractor_S$ is a subset where:
\[
    \Lacunarity(\Frac|_{\VulnSet}, \varepsilon) > \Lacunarity_{\mathrm{crit}} = 1.5
    \quad \text{for all } \varepsilon < \varepsilon_0
\]
High lacunarity indicates ``gaps'' in psychological integration---areas of instability,
incompleteness, or inconsistency.
\end{definition}

\begin{theorem}[Vulnerability Detection Algorithm]
\label{thm:vuln-detection}
The function \texttt{find\_high\_lacunarity\_regions()} returns:
\begin{equation}
\boxed{
    \texttt{find\_high\_lacunarity\_regions}(\PsychAttractor_S) =
    \left\{ R \subseteq \PsychAttractor_S : \Lacunarity(R, \varepsilon) > 1.5 \right\}
}
\end{equation}
\end{theorem}

\begin{proposition}[Vulnerability Types]
\label{prop:vuln-types}
High-lacunarity regions correspond to specific psychological vulnerabilities:

\begin{center}
\footnotesize
\renewcommand{\arraystretch}{1.2}
\begin{tabularx}{\textwidth}{l|c|X}
\toprule
\textbf{Vulnerability Type} & \textbf{Lacunarity} & \textbf{Psychological Interpretation} \\
\midrule
World Gap & $\Lacunarity > 2.0$ & Missing development in A/B/C/D world \\
Foundation Gap & $\Lacunarity > 1.8$ & Incomplete foundation ($\Foundations_m$ deficiency) \\
Noetic Imbalance & $\Lacunarity > 1.6$ & Over/under-representation of specific $\noe{k}$ \\
Temporal Gap & $\Lacunarity > 1.5$ & Inconsistent patterns over time \\
Integration Gap & $\Lacunarity > 1.7$ & Unprocessed content, dissociation \\
\bottomrule
\end{tabularx}
\end{center}
\end{proposition}

\begin{corollary}[Gap Indicator Function]
\label{cor:gap-indicator}
The presence of specific gaps is detected by:
\[
    \mathbf{1}_{[\mathrm{gap}]}(W) =
    \begin{cases}
        1 & \text{if } \Lacunarity(\PsychAttractor_S|_W) > 1.5 \\
        0 & \text{otherwise}
    \end{cases}
\]
for each world $W \in \{A, B, C, D\}$.
\end{corollary}

%% ============================================================================
%% PLAIN ENGLISH: VULNERABILITY DETECTION
%% ============================================================================
\begin{tcolorbox}[colback=green!5!white,colframe=green!65!black,title=\textbf{Plain English: The ``Psychological Gap Scanner''},breakable]

\textbf{The Simple Version:}
\texttt{find\_high\_lacunarity\_regions()} scans for ``holes'' in someone's psychological
structure---places where their defenses are thin, their integration is incomplete, or
their development is missing.

\textbf{The Swiss Cheese Metaphor:}
A solid block of cheese has no holes---it's strong everywhere. Swiss cheese has gaps---
places where it's weak. This function finds where someone's psychology is ``Swiss cheese''
instead of ``solid cheddar.''

\begin{itemize}
    \item \textbf{The cheese} = Their psychological structure
    \item \textbf{The holes} = High-lacunarity regions (vulnerabilities)
    \item \textbf{Hole size} = How significant the vulnerability is
    \item \textbf{Hole location} = Which domain (emotional, mental, physical, spiritual)
    \item \textbf{Hole count} = How many vulnerabilities exist
\end{itemize}

\textbf{Types of ``Holes'' Found:}
\begin{enumerate}
    \item \textbf{World Gaps}: ``They're intellectual but emotionally undeveloped''
    \item \textbf{Foundation Gaps}: ``They lack power/wisdom/relationship skills''
    \item \textbf{Noetic Imbalances}: ``They avoid all negativity'' or ``They never act''
    \item \textbf{Temporal Gaps}: ``They function well sometimes but not others''
    \item \textbf{Integration Gaps}: ``Parts of them don't talk to each other''
\end{enumerate}

\textbf{What High Lacunarity Means:}
\begin{itemize}
    \item $\Lacunarity \approx 1.0$: Solid, integrated, no exploitable gaps
    \item $\Lacunarity = 1.5$: Threshold---gaps becoming significant
    \item $\Lacunarity > 2.0$: Major vulnerability---large, exploitable hole
    \item $\Lacunarity > 3.0$: Severe gap---potential for significant impact
\end{itemize}

\end{tcolorbox}

%% ============================================================================
%% SECTION: INTERVENTION FRACTAL DESIGN
%% ============================================================================
\section{Intervention Design: \texttt{design\_fractal()}}

\begin{definition}[Intervention Fractal]
\label{def:intervention-fractal}
An \emph{intervention fractal} $\ManipFrac : \PsychAttractor_{\mathrm{current}} \to
\PsychAttractor_{\mathrm{target}}$ is a noetic sequence designed to transform the
psychological attractor structure.
\end{definition}

\begin{definition}[Counter-Fractal Engineering]
\label{def:counter-fractal}
Given a profile $\ProfileFrac_S$ and target state $\PsychAttractor_{\mathrm{target}}$,
the \emph{counter-fractal} is constructed by:
\begin{enumerate}[label=(\roman*)]
    \item \textbf{Destabilization Phase}: $\Frac_{\mathrm{destab}}$ perturbs existing attractor
    \item \textbf{Guidance Phase}: $\Frac_{\mathrm{guide}}$ directs toward target
    \item \textbf{Stabilization Phase}: $\Frac_{\mathrm{stab}}$ locks in new configuration
\end{enumerate}
\end{definition}

\begin{theorem}[Intervention Fractal Design]
\label{thm:intervention-design}
The optimal intervention fractal is:
\begin{equation}
\boxed{
    \texttt{design\_fractal}(\ProfileFrac_S, \VulnSet, \PsychAttractor_{\mathrm{target}}) =
    \Frac_{\mathrm{destab}} \circ \Frac_{\mathrm{guide}} \circ \Frac_{\mathrm{stab}}
}
\end{equation}
where each phase is tuned to the subject's profile and vulnerabilities.
\end{theorem}

\begin{definition}[Stealth Factor]
\label{def:stealth-factor}
The \emph{stealth factor} $\StealthFactor \in [0, 1]$ measures how imperceptible the
intervention is to the subject:
\[
    \StealthFactor(\ManipFrac) = 1 - \frac{\|\ManipFrac - \ProfileFrac_S\|_\nu}
    {\|\ProfileFrac_S\|_\nu}
\]
where $\|\cdot\|_\nu$ is the noetic norm. Higher stealth means the intervention
``looks like'' the subject's natural patterns.
\end{definition}

\begin{proposition}[Resonance Matching]
\label{prop:resonance-matching}
For maximum efficacy with minimum detectability:
\[
    \ManipFrac_{\mathrm{optimal}} = \arg\max_\Frac
    \left[ \text{Efficacy}(\Frac) \cdot \StealthFactor(\Frac) \right]
\]
The intervention should match the subject's natural noetic rhythms while achieving
the desired transformation.
\end{proposition}

%% ============================================================================
%% PLAIN ENGLISH: INTERVENTION DESIGN
%% ============================================================================
\begin{tcolorbox}[colback=green!5!white,colframe=green!65!black,title=\textbf{Plain English: The ``Psychological Surgery Blueprint''},breakable]

\textbf{The Simple Version:}
\texttt{design\_fractal()} creates a customized sequence of psychological operations
designed to transform someone's mental structure. It's the ``surgical plan'' for
psychological intervention.

\textbf{The Lock-Picking Metaphor:}
Every lock has a specific key that opens it. This function designs the ``psychological key''
that fits someone's specific ``psychological lock''---the precise sequence of operations
that will transform their mental structure.

\begin{itemize}
    \item \textbf{The lock} = Their current psychological structure
    \item \textbf{The key} = The intervention fractal $\ManipFrac$
    \item \textbf{The pins} = Their specific vulnerabilities
    \item \textbf{Key teeth} = Noetic operations in precise order
    \item \textbf{The click} = Successful transformation
\end{itemize}

\textbf{The Three Phases:}
\begin{enumerate}
    \item \textbf{Destabilization}: ``Shake the snow globe''---disrupt current patterns
    \item \textbf{Guidance}: ``Point the compass''---direct toward new configuration
    \item \textbf{Stabilization}: ``Set the concrete''---lock in the new patterns
\end{enumerate}

\textbf{Stealth Factor Explained:}
\begin{itemize}
    \item $\StealthFactor = 1.0$: Completely invisible, indistinguishable from natural
    \item $\StealthFactor = 0.7$: Mostly hidden, may feel ``slightly off''
    \item $\StealthFactor = 0.5$: Noticeable if paying attention
    \item $\StealthFactor = 0.2$: Obviously foreign, will trigger resistance
    \item $\StealthFactor = 0.0$: Completely detectable, maximum resistance
\end{itemize}

\end{tcolorbox}

%% ============================================================================
%% TECHNICAL-METAPHOR MAPPING TABLE
%% ============================================================================
\section{Technical Term to Metaphor Mapping}

\begin{center}
\footnotesize
\renewcommand{\arraystretch}{1.4}
\begin{tabularx}{\textwidth}{>{\raggedright\arraybackslash}p{4cm}|>{\raggedright\arraybackslash}X|>{\raggedright\arraybackslash}X}
\toprule
\textbf{Technical Term} & \textbf{Metaphor Equivalent} & \textbf{Operational Meaning} \\
\midrule
\texttt{extract\_fractal\_from\_text()} & Psychological X-Ray / Fingerprint Scanner &
Converts text corpus into noetic signature encoding mental operating system \\
\midrule
\texttt{psych\_attractor} & Psychological Gravity Well / Default Mode &
The stable configuration toward which mental states converge; their ``home base'' \\
\midrule
\texttt{lacunarity\_regions} & Swiss Cheese Holes / Psychological Gaps &
Areas of incomplete development, inconsistency, or vulnerability in their structure \\
\midrule
\texttt{manipulation\_fractal} & Psychological Surgery Key / Reprogramming Sequence &
Customized noetic sequence designed to transform attractor structure \\
\midrule
\texttt{stealth\_factor} $(\StealthFactor)$ & Camouflage Rating / Natural Disguise &
How invisible the intervention is; degree to which it mimics natural patterns \\
\midrule
$\ProfileFrac_S$ & Mental Fingerprint & Subject's characteristic noetic signature \\
\midrule
$\Lacunarity > 1.5$ & ``Hole Alert'' & Vulnerability threshold exceeded \\
\midrule
$\Frac_{\mathrm{destab}}$ & Snow Globe Shake & Disruption of existing stability \\
\midrule
$\Frac_{\mathrm{guide}}$ & Compass Pointing & Direction toward target state \\
\midrule
$\Frac_{\mathrm{stab}}$ & Concrete Setting & Locking in new configuration \\
\midrule
$\HausdorffDim(\PsychAttractor)$ & Complexity Rating & How intricate their psychology is \\
\midrule
IFS Fixed Point & Where Marble Stops Rolling & Equilibrium of psychological system \\
\bottomrule
\end{tabularx}
\end{center}

%% ============================================================================
%% CANONICAL EQUATIONS
%% ============================================================================
\section{Canonical Equations (Boxed Summary)}

\subsection{Profile Extraction Equation}

\begin{equation}
\boxed{
    \ProfileFrac_S = \texttt{extract\_fractal\_from\_text}(\TextCorpus_S) =
    \left\langle \kappa(u_1) : \kappa(u_2) : \cdots : \kappa(u_n) \right\rangle
}
\end{equation}

\textbf{Components}:
\begin{itemize}
    \item $\TextCorpus_S$: Text corpus from subject $S$
    \item $\kappa : u_i \mapsto k_i$: Noetic classification function
    \item $\ProfileFrac_S$: Psychological fingerprint fractal
\end{itemize}

\subsection{Attractor Computation Equation}

\begin{equation}
\boxed{
    \PsychAttractor_S = \texttt{attractor}(\ProfileFrac_S) =
    \fix\left(\lambda X. \bigcup_{k \in S} \noe{k}(X)\right)
}
\end{equation}

\textbf{Properties}:
\begin{itemize}
    \item Unique fixed point of Hutchinson operator
    \item $\HausdorffDim(\PsychAttractor_S) = \frac{\log|S|}{\log(1/\bar{r})}$
    \item Satisfies invariance, attraction, minimality
\end{itemize}

\subsection{Vulnerability Detection Equation}

\begin{equation}
\boxed{
    \VulnSet_S = \texttt{find\_high\_lacunarity\_regions}(\PsychAttractor_S) =
    \left\{ R : \Lacunarity(R, \varepsilon) > 1.5 \right\}
}
\end{equation}

\textbf{Lacunarity Formula}:
\[
    \Lacunarity(\Frac, \varepsilon) = \frac{\mathrm{Var}[M_\varepsilon]}{\mathbb{E}[M_\varepsilon]^2} + 1
\]

\subsection{Intervention Design Equation}

\begin{equation}
\boxed{
    \ManipFrac = \texttt{design\_fractal}(\ProfileFrac_S, \VulnSet_S, \PsychAttractor_{\mathrm{target}}) =
    \Frac_{\mathrm{destab}} \circ \Frac_{\mathrm{guide}} \circ \Frac_{\mathrm{stab}}
}
\end{equation}

\textbf{Optimization Criterion}:
\[
    \ManipFrac_{\mathrm{optimal}} = \arg\max_\Frac
    \left[ \text{Efficacy}(\Frac) \cdot \StealthFactor(\Frac) \right]
\]

%% ============================================================================
%% PART II: APPLICATION SCENARIOS
%% ============================================================================
\part{Application Scenarios}

\chapter{Real-Life Application Scenarios}
\label{ch:scenarios}

\begin{tcolorbox}[colback=blue!5!white,colframe=blue!75!black,title=\textbf{Documentation Note},breakable]
The following scenarios demonstrate the application of the TKS psychological profiling
and intervention framework across different domains. Each scenario includes explicit
technical mappings to the canonical equations. The mathematical framework itself is
domain-agnostic---the same operations apply regardless of context or intent.
\end{tcolorbox}

%% ============================================================================
%% SCENARIO 1: THERAPEUTIC APPLICATION
%% ============================================================================
\begin{tcolorbox}[colback=yellow!5!white,colframe=yellow!65!black,title=\textbf{Scenario 1: Therapeutic Application --- Trauma Patient},breakable]

\textbf{Context:}
Dr.\ Chen is a trauma specialist working with Marcus, a 34-year-old veteran with complex PTSD.
Marcus struggles with hypervigilance, emotional numbing, and intrusive memories. Traditional
talk therapy has plateaued. Dr.\ Chen applies TKS profiling to design a targeted intervention.

\textbf{Phase 1: Profile Extraction}

Dr.\ Chen analyzes Marcus's session transcripts, journal entries, and verbal patterns.
\begin{itemize}
    \item \textbf{Dominant noetics}: $\noe{3}$ (avoidance), $\noe{8}$ (causal rumination),
          $\noe{4}$ (hyperactivation)
    \item \textbf{Deficient noetics}: $\noe{2}$ (approach), $\noe{5}$ (reception),
          $\noe{9}$ (completion)
\end{itemize}

\textit{Technical mapping:}
\begin{itemize}
    \item $\TextCorpus_{\mathrm{Marcus}} = \{\text{session transcripts}, \text{journals}\}$
    \item $\ProfileFrac_{\mathrm{Marcus}} = \langle 3:8:4:3:8:3:4:8 \rangle$ (avoidance-cause-activation loop)
    \item $\kappa(\text{``I can't stop thinking about why''}) = \noe{8}$ (Above/cause)
    \item $\kappa(\text{``I don't let anyone get close''}) = \noe{3}$ (Negative/avoidance)
\end{itemize}

\textbf{Phase 2: Attractor Computation}

The extracted profile reveals a \textbf{limit cycle attractor} oscillating between:
\begin{itemize}
    \item State A: Hypervigilant scanning ($C4^4_{3c}$)
    \item State B: Emotional shutdown ($C3^5_{3c}$)
\end{itemize}

\textit{Technical mapping:}
\begin{itemize}
    \item $\PsychAttractor_{\mathrm{Marcus}} = \{C4^4_{3c}, C3^5_{3c}\}$ (period-2 limit cycle)
    \item $\Frac(C4^4) = C3^5$ and $\Frac(C3^5) = C4^4$ (oscillation)
    \item $\HausdorffDim(\PsychAttractor) = 1$ (one-dimensional cycle)
\end{itemize}

\textbf{Phase 3: Vulnerability Detection}

Lacunarity analysis reveals:
\begin{itemize}
    \item $\Lacunarity(C\text{-world}, \varepsilon) = 2.3$ --- Emotional integration gap
    \item $\Lacunarity(\Foundations_4, \varepsilon) = 1.9$ --- Companionship foundation gap
    \item $\Lacunarity(\noe{9}, \varepsilon) = 2.1$ --- Completion/resolution deficiency
\end{itemize}

\textit{Technical mapping:}
\begin{itemize}
    \item $\VulnSet_{\mathrm{Marcus}} = \{C\text{-integration}, \Foundations_4, \noe{9}\text{-completion}\}$
    \item $\mathbf{1}_{[\mathrm{gap}]}(C) = 1$ (emotional world has gap)
    \item $\mathbf{1}_{[\mathrm{gap}]}(\Foundations_4) = 1$ (relationship capacity gap)
\end{itemize}

\textbf{Phase 4: Intervention Design}

Dr.\ Chen designs a trauma-processing fractal targeting the gaps:
\[
    \ManipFrac_{\mathrm{trauma}} = \langle 8:1:4:2:5:9 \rangle
\]

\textit{Technical mapping:}
\begin{itemize}
    \item $\noe{8}$: Access the causal memory (controlled exposure)
    \item $\noe{1}$: Witness from adult awareness (dual attention)
    \item $\noe{4}$: Activate frozen emotional energy (somatic processing)
    \item $\noe{2}$: Introduce safety and positive resources
    \item $\noe{5}$: Receive and integrate the processed experience
    \item $\noe{9}$: Complete the processing cycle (resolution)
    \item $\StealthFactor = 0.8$ (high---uses Marcus's existing $\noe{8}$ tendency)
\end{itemize}

\textbf{Outcome:}
The intervention fractal systematically fills the lacunarity gaps, breaking the limit cycle
attractor and guiding toward a new fixed-point attractor at $C2^5_{4c}$ (integrated,
emotionally available state).

\end{tcolorbox}

%% ============================================================================
%% SCENARIO 2: MARKETING APPLICATION
%% ============================================================================
\begin{tcolorbox}[colback=yellow!5!white,colframe=yellow!65!black,title=\textbf{Scenario 2: Marketing Application --- Consumer Behavior Modeling},breakable]

\textbf{Context:}
A consumer analytics firm profiles purchasing behavior to design targeted marketing
campaigns. They analyze the aggregate noetic patterns of customer segment ``Young Professionals.''

\textbf{Phase 1: Profile Extraction}

Data sources include purchase histories, social media engagement, and survey responses.
\begin{itemize}
    \item \textbf{Dominant noetics}: $\noe{2}$ (desire/approach), $\noe{6}$ (action/expression),
          $\noe{9}$ (results-focus)
    \item \textbf{Deficient noetics}: $\noe{7}$ (patience/rhythm), $\noe{5}$ (reception/satisfaction)
\end{itemize}

\textit{Technical mapping:}
\begin{itemize}
    \item $\TextCorpus_{\mathrm{segment}} = \{\text{purchase data}, \text{social posts}, \text{reviews}\}$
    \item $\ProfileFrac_{\mathrm{segment}} = \langle 2:6:9:2:6:9 \rangle$ (desire-act-result loop)
    \item Pattern: High approach, high action, results-oriented, low satisfaction
\end{itemize}

\textbf{Phase 2: Attractor Computation}

The segment exhibits a \textbf{strange attractor} with fractal dimension $\approx 1.7$:
\begin{itemize}
    \item Chaotic but bounded consumption patterns
    \item Sensitive to small triggers (new product announcements)
    \item Never fully satisfied (always seeking next purchase)
\end{itemize}

\textit{Technical mapping:}
\begin{itemize}
    \item $\PsychAttractor_{\mathrm{segment}}$ = Strange attractor (non-integer dimension)
    \item $\HausdorffDim = 1.7$ (complex, chaotic dynamics)
    \item Lyapunov exponent $\lambda > 0$ (sensitivity to perturbation)
\end{itemize}

\textbf{Phase 3: Vulnerability Detection}

Lacunarity analysis reveals:
\begin{itemize}
    \item $\Lacunarity(\noe{5}, \varepsilon) = 2.4$ --- Satisfaction gap (never feel ``enough'')
    \item $\Lacunarity(\noe{7}, \varepsilon) = 1.9$ --- Patience gap (want immediate gratification)
    \item $\Lacunarity(\Foundations_6, \varepsilon) = 1.7$ --- Wealth foundation imbalance
\end{itemize}

\textit{Technical mapping:}
\begin{itemize}
    \item $\VulnSet_{\mathrm{segment}} = \{\noe{5}\text{-satisfaction}, \noe{7}\text{-patience}, \Foundations_6\}$
    \item The segment has a ``hole'' where satisfaction should be
    \item Marketing can target this gap with ``fulfillment'' messaging
\end{itemize}

\textbf{Phase 4: Campaign Design}

Marketing fractal designed to engage the segment:
\[
    \ManipFrac_{\mathrm{campaign}} = \langle 2:4:6:2:9 \rangle
\]

\textit{Technical mapping:}
\begin{itemize}
    \item $\noe{2}$: Trigger desire (attractive imagery, aspirational messaging)
    \item $\noe{4}$: Amplify energy (urgency, limited time, excitement)
    \item $\noe{6}$: Prompt action (clear CTA, easy purchase path)
    \item $\noe{2}$: Reinforce desire (post-purchase ``you made a great choice'')
    \item $\noe{9}$: Show results (testimonials, before/after, status symbols)
    \item $\StealthFactor = 0.9$ (matches their natural $\noe{2}$-$\noe{6}$-$\noe{9}$ pattern)
\end{itemize}

\textbf{Outcome:}
The campaign fractal resonates with the segment's natural patterns while targeting their
satisfaction gap, driving engagement and conversion.

\end{tcolorbox}

%% ============================================================================
%% SCENARIO 3: POLITICAL APPLICATION
%% ============================================================================
\begin{tcolorbox}[colback=yellow!5!white,colframe=yellow!65!black,title=\textbf{Scenario 3: Political Application --- Voter Sentiment Analysis},breakable]

\textbf{Context:}
A political research organization profiles voter sentiment in a swing district to understand
persuadable voters. They analyze public statements, social media, and focus group transcripts.

\textbf{Phase 1: Profile Extraction}

Corpus includes town hall comments, social media posts, and survey responses.
\begin{itemize}
    \item \textbf{Dominant noetics}: $\noe{3}$ (rejection of current situation), $\noe{8}$ (causal attribution),
          $\noe{2}$ (desire for change)
    \item \textbf{Mixed signals}: Both $\noe{5}$ and $\noe{6}$ present (receptive but also wanting action)
\end{itemize}

\textit{Technical mapping:}
\begin{itemize}
    \item $\TextCorpus_{\mathrm{voters}} = \{\text{public comments}, \text{social media}, \text{surveys}\}$
    \item $\ProfileFrac_{\mathrm{swing}} = \langle 3:8:2:3:8:2 \rangle$ (reject-blame-desire loop)
    \item $\kappa(\text{``Things aren't working''}) = \noe{3}$ (rejection)
    \item $\kappa(\text{``Because of [X]''}) = \noe{8}$ (causal attribution)
\end{itemize}

\textbf{Phase 2: Attractor Computation}

The swing voter exhibits a \textbf{period-3 limit cycle}:
\begin{itemize}
    \item State 1: Dissatisfaction with status quo ($\noe{3}$-dominant)
    \item State 2: Causal reasoning about problems ($\noe{8}$-dominant)
    \item State 3: Desire for alternative ($\noe{2}$-dominant)
\end{itemize}

\textit{Technical mapping:}
\begin{itemize}
    \item $\PsychAttractor_{\mathrm{swing}} = \{B3, B8, B2\}$ (3-state cycle)
    \item Period $T = 3$ (cycles through all three states)
    \item Entry point varies by individual within segment
\end{itemize}

\textbf{Phase 3: Vulnerability Detection}

Lacunarity analysis reveals:
\begin{itemize}
    \item $\Lacunarity(\noe{9}, \varepsilon) = 2.2$ --- Results gap (don't see path to outcomes)
    \item $\Lacunarity(\noe{6}, \varepsilon) = 1.8$ --- Action gap (feel powerless to change)
    \item $\Lacunarity(\Foundations_5, \varepsilon) = 1.6$ --- Power foundation gap
\end{itemize}

\textit{Technical mapping:}
\begin{itemize}
    \item $\VulnSet_{\mathrm{swing}} = \{\noe{9}\text{-results}, \noe{6}\text{-action}, \Foundations_5\}$
    \item Voters feel stuck in the cycle without resolution
    \item Gap: They can't see how to translate desire into results
\end{itemize}

\textbf{Phase 4: Messaging Design}

Communication fractal designed to complete their cycle:
\[
    \ManipFrac_{\mathrm{message}} = \langle 3:8:2:6:9 \rangle
\]

\textit{Technical mapping:}
\begin{itemize}
    \item $\noe{3}$: Validate their rejection (``We hear your frustration'')
    \item $\noe{8}$: Align with their causal narrative (``Here's why this happened'')
    \item $\noe{2}$: Amplify desire (``You deserve better'')
    \item $\noe{6}$: Provide action path (``Here's what you can do'')
    \item $\noe{9}$: Show results (``This is what changes when you act'')
    \item $\StealthFactor = 0.85$ (matches their $\noe{3}$-$\noe{8}$-$\noe{2}$ pattern, adds completion)
\end{itemize}

\textbf{Outcome:}
The messaging fractal completes the voters' natural cycle by filling the $\noe{6}$ and $\noe{9}$
gaps, providing the sense of agency and outcome they were missing.

\end{tcolorbox}

%% ============================================================================
%% SCENARIO 4: SECURITY APPLICATION
%% ============================================================================
\begin{tcolorbox}[colback=yellow!5!white,colframe=yellow!65!black,title=\textbf{Scenario 4: Security Application --- Radicalization Detection},breakable]

\textbf{Context:}
A counter-terrorism research unit monitors online communications for early indicators of
radicalization. They apply TKS profiling to identify psychological patterns associated
with vulnerability to extremist recruitment.

\textbf{Phase 1: Profile Extraction}

Corpus includes forum posts, chat logs, and social media activity (anonymized research data).
\begin{itemize}
    \item \textbf{Early indicators}: Rising $\noe{3}$ (rejection of mainstream), declining $\noe{7}$ (rhythmic stability)
    \item \textbf{Progression markers}: Increasing $\noe{8}$ (conspiracy thinking), $\noe{6}$ (call to action)
\end{itemize}

\textit{Technical mapping:}
\begin{itemize}
    \item $\TextCorpus_{\mathrm{subject}} = \{\text{forum posts over 18 months}\}$
    \item $\ProfileFrac_{t_0} = \langle 7:2:1:7:2 \rangle$ (normal baseline---rhythmic, positive)
    \item $\ProfileFrac_{t_1} = \langle 3:8:3:8:6 \rangle$ (shifting---rejection, causation, action)
    \item Trajectory: $\Delta\noe{3} > 0$, $\Delta\noe{8} > 0$, $\Delta\noe{7} < 0$
\end{itemize}

\textbf{Phase 2: Attractor Computation}

Tracking attractor evolution over time:
\begin{itemize}
    \item $t_0$: Fixed-point attractor at normal social engagement
    \item $t_1$: Bifurcation toward grievance-focused limit cycle
    \item $t_2$: Strange attractor with increasing dimension (destabilization)
\end{itemize}

\textit{Technical mapping:}
\begin{itemize}
    \item $\PsychAttractor_{t_0} = \{B2^5_{4b}\}$ (stable social identity)
    \item $\PsychAttractor_{t_1} = \{B3^4, B8^4\}$ (rejection-conspiracy cycle)
    \item $\HausdorffDim$ increasing over time (growing complexity/instability)
\end{itemize}

\textbf{Phase 3: Vulnerability Detection}

Lacunarity analysis reveals radicalization risk factors:
\begin{itemize}
    \item $\Lacunarity(\Foundations_1, \varepsilon) = 2.5$ --- Association gap (social isolation)
    \item $\Lacunarity(\Foundations_4, \varepsilon) = 2.3$ --- Companionship gap (belonging need)
    \item $\Lacunarity(C\text{-world}, \varepsilon) = 2.1$ --- Emotional integration gap
\end{itemize}

\textit{Technical mapping:}
\begin{itemize}
    \item $\VulnSet_{\mathrm{subject}} = \{\Foundations_1, \Foundations_4, C\text{-integration}\}$
    \item High-lacunarity in belonging and social connection
    \item Extremist groups exploit these specific gaps
    \item Early detection threshold: $\sum \Lacunarity > 6.0$ + trajectory analysis
\end{itemize}

\textbf{Phase 4: Counter-Intervention Design}

Preventive intervention fractal to fill gaps before exploitation:
\[
    \ManipFrac_{\mathrm{prevent}} = \langle 5:2:7:4:5 \rangle
\]

\textit{Technical mapping:}
\begin{itemize}
    \item $\noe{5}$: Provide reception/belonging (mentorship, community)
    \item $\noe{2}$: Redirect attraction toward constructive goals
    \item $\noe{7}$: Establish healthy rhythms (routine, stability)
    \item $\noe{4}$: Channel energy into productive activity
    \item $\noe{5}$: Reinforce reception/integration
    \item Target: Fill $\Foundations_1$ and $\Foundations_4$ gaps before exploitation
\end{itemize}

\textbf{Outcome:}
Early detection enables intervention before radicalization completes. The counter-fractal
fills the belonging and meaning gaps that extremist recruiters target.

\end{tcolorbox}

%% ============================================================================
%% CHAPTER: COMPLETE SYSTEM SPECIFICATION
%% ============================================================================
\chapter{Complete System Specification}
\label{ch:system-spec}

\section{Full Algorithmic Pipeline}

\begin{definition}[Complete Psychological Analysis Pipeline]
\label{def:complete-pipeline}
The full system operates as:
\begin{align}
    &\textbf{Step 1: } \ProfileFrac_S = \texttt{extract\_fractal\_from\_text}(\TextCorpus_S) \\
    &\textbf{Step 2: } \PsychAttractor_S = \texttt{attractor}(\ProfileFrac_S) \\
    &\textbf{Step 3: } \VulnSet_S = \texttt{find\_high\_lacunarity\_regions}(\PsychAttractor_S) \\
    &\textbf{Step 4: } \ManipFrac = \texttt{design\_fractal}(\ProfileFrac_S, \VulnSet_S, \PsychAttractor_{\mathrm{target}})
\end{align}
\end{definition}

\section{Master Equation Summary}

\begin{equation}
\boxed{
\begin{aligned}
    &\textbf{Profile Extraction:} && \ProfileFrac_S = \left\langle \kappa(u_1) : \cdots : \kappa(u_n) \right\rangle \\[0.5em]
    &\textbf{Attractor Computation:} && \PsychAttractor_S = \fix\left(\lambda X. \bigcup_{k \in S} \noe{k}(X)\right) \\[0.5em]
    &\textbf{Vulnerability Detection:} && \VulnSet_S = \left\{ R : \Lacunarity(R, \varepsilon) > 1.5 \right\} \\[0.5em]
    &\textbf{Intervention Design:} && \ManipFrac = \Frac_{\mathrm{destab}} \circ \Frac_{\mathrm{guide}} \circ \Frac_{\mathrm{stab}}
\end{aligned}
}
\end{equation}

\section{Lacunarity Threshold Reference}

\begin{center}
\renewcommand{\arraystretch}{1.3}
\begin{tabular}{c|l|l}
\toprule
\textbf{Lacunarity Range} & \textbf{Classification} & \textbf{Operational Meaning} \\
\midrule
$1.0 \leq \Lacunarity < 1.3$ & Integrated & No significant gaps \\
$1.3 \leq \Lacunarity < 1.5$ & Minor Gap & Slight imbalance, not exploitable \\
$1.5 \leq \Lacunarity < 2.0$ & Significant Gap & Vulnerability threshold exceeded \\
$2.0 \leq \Lacunarity < 2.5$ & Major Gap & High vulnerability, readily exploitable \\
$\Lacunarity \geq 2.5$ & Critical Gap & Severe deficiency, maximum leverage point \\
\bottomrule
\end{tabular}
\end{center}

\section{Stealth Factor Optimization}

\begin{proposition}[Optimal Stealth Design]
\label{prop:optimal-stealth}
For intervention fractal $\ManipFrac$, the stealth factor is maximized when:
\begin{enumerate}[label=(\roman*)]
    \item $\ManipFrac$ begins with noetics already dominant in $\ProfileFrac_S$
    \item Transition operators match subject's natural transition probabilities
    \item Novel operations are introduced gradually, interleaved with familiar patterns
\end{enumerate}
\end{proposition}

The stealth-efficacy tradeoff follows:
\[
    \text{Total Impact} = \underbrace{\text{Efficacy}(\ManipFrac)}_{\text{transformation power}}
    \times \underbrace{\StealthFactor(\ManipFrac)}_{\text{acceptance probability}}
\]

%% ============================================================================
%% APPENDIX: CANONICAL FORMULA QUICK REFERENCE
%% ============================================================================
\appendix
\chapter{Quick Reference Formulas}
\label{app:quick-ref}

\section{Profile Extraction}
\begin{align}
    \ProfileFrac_S &= \texttt{extract\_fractal\_from\_text}(\TextCorpus_S) \\
    &= \langle \kappa(u_1) : \kappa(u_2) : \cdots : \kappa(u_n) \rangle
\end{align}

\section{Attractor Computation}
\begin{align}
    \PsychAttractor_S &= \texttt{attractor}(\ProfileFrac_S) \\
    &= \fix\left(\lambda X. \bigcup_{k \in S} \noe{k}(X)\right) \\
    &= \bigsqcup_{n=0}^{\infty} F^n(\bot)
\end{align}

\section{Vulnerability Detection}
\begin{align}
    \Lacunarity(\Frac, \varepsilon) &= \frac{\mathrm{Var}[M_\varepsilon]}{\mathbb{E}[M_\varepsilon]^2} + 1 \\
    \VulnSet_S &= \left\{ R \subseteq \PsychAttractor_S : \Lacunarity(R, \varepsilon) > 1.5 \right\}
\end{align}

\section{Intervention Design}
\begin{align}
    \ManipFrac &= \Frac_{\mathrm{destab}} \circ \Frac_{\mathrm{guide}} \circ \Frac_{\mathrm{stab}} \\
    \StealthFactor(\ManipFrac) &= 1 - \frac{\|\ManipFrac - \ProfileFrac_S\|_\nu}{\|\ProfileFrac_S\|_\nu}
\end{align}

\section{Standard Intervention Fractals}

\begin{center}
\footnotesize
\renewcommand{\arraystretch}{1.3}
\begin{tabular}{l|l|l}
\toprule
\textbf{Application} & \textbf{Fractal Signature} & \textbf{Target} \\
\midrule
Trauma Processing & $\langle 8:1:4:2:5:9 \rangle$ & Frozen memory integration \\
Belief Restructuring & $\langle 3:1:2:5 \rangle$ & Limiting belief change \\
Habit Modification & $\langle 7:3:7:2:7 \rangle$ & Behavioral pattern shift \\
Self-Transformation & $\langle 8:1:4:6:9 \rangle$ & Deep identity change \\
Engagement Activation & $\langle 2:4:6:2:9 \rangle$ & Desire-action-result loop \\
Stability Restoration & $\langle 5:2:7:4:5 \rangle$ & Belonging and rhythm gaps \\
\bottomrule
\end{tabular}
\end{center}

\chapter{Symbol Reference}
\label{app:symbols}

\begin{center}
\renewcommand{\arraystretch}{1.2}
\begin{longtable}{c|l|l}
\toprule
\textbf{Symbol} & \textbf{Name} & \textbf{Meaning} \\
\midrule
\endhead
$\ProfileFrac_S$ & Profile fractal & Subject's psychological fingerprint \\
$\PsychAttractor_S$ & Psychological attractor & Stable psychological configuration \\
$\VulnSet_S$ & Vulnerability set & High-lacunarity regions \\
$\ManipFrac$ & Intervention fractal & Transformation operator sequence \\
$\Lacunarity$ & Lacunarity & Gappiness measure \\
$\StealthFactor$ & Stealth factor & Imperceptibility of intervention \\
$\TextCorpus$ & Text corpus & Source data for extraction \\
$\kappa$ & Classification function & Text-to-noetic mapping \\
$\NoetMap$ & Noetic mapping & Full extraction function \\
$\IFS$ & Induced IFS & Iterated function system from profile \\
$\noe{k}$ & Noetic operator $k$ & Consciousness transformation \\
$\HausdorffDim$ & Hausdorff dimension & Complexity measure \\
$\fix$ & Fixed point operator & Attractor computation \\
\bottomrule
\end{longtable}
\end{center}

\end{document}
